\clearpage
\part{Correspondence and the next mathematical steps}
\section{What has been connected}
The construction now has two complete mathematical objects on the same physical space. One is the unitary evolution obtained by assembling exact mode exponentials and conjugating through Fourier synthesis. The other is the self-adjoint unbounded Hamiltonian obtained from the maximal coefficient multiplier and its exact first-order domain.

The theorem chain can be summarized without collapsing its distinct obligations.
\begin{longtable}{@{}L{.13\linewidth}L{.35\linewidth}L{.43\linewidth}@{}}
\toprule
Blocks & Construction & Obligation discharged\\\midrule\endhead
41--42 & Mode symbols and finite torus packets & Exact mass shell and generator consistency in a physical packet norm.\\
43--44 & Exact mode exponentials in algebraic and Hilbert realizations & Group identities, Hilbert unitarity, and packet norm preservation.\\
45 & Compact time windows & One finite-packet consistency bound valid throughout a fixed interval.\\
46 & Completed coefficient space & Full unitary maps and cutoff removal for every square-summable coefficient field.\\
47 & Physical Fourier equivalence & Surjective synthesis and lossless transport of evolution and convergence to torus fields.\\
48 & Maximal coefficient operator & Density, adjoint-domain equality, closedness, and simultaneous state/image approximation.\\
49 & Weighted and physical domains & Exact first-order characterization, mass independence, and physical self-adjoint Hamiltonian.\\
\bottomrule
\end{longtable}

The source records 147 active mathematical modules, 3,920 named declarations, 2,968 theorem declarations, and 207 negative controls at this checkpoint. Block 49 adds 40 named declarations, including 30 theorems. These are versioned receipt totals. Supporting lemmas contribute to them, so they are not counts of independent physical discoveries.

The recorded Block 49 checks include warning-as-error compilation, the 3,064-job module import closure, checks of all 40 public declarations, theorem-level axiom inspection, and rejection of three deliberately false controls. The permitted imported logical dependencies observed in that audit are \path{Classical.choice}, \path{Quot.sound}, and \path{propext}. The new-source scan records no \path{sorry}, \path{admit}, custom axiom, \path{native_decide}, or \path{unsafe}.

Production, fresh-extraction, and final-archive replay records are included with the source handoff. This monograph preparation inspected those records without rerunning Lean. The pinned Lean toolchain and Mathlib dependencies were not rebuilt from their own sources by the recorded checkpoint procedure.

\section{Ordinary deductions and checked correspondences}
The finite-packet exposition preserves several useful deductions that go beyond a direct one-to-one declaration map. Their proofs remain part of the mathematics of the volume. Their status is stated here so that the formal record can be read precisely.

\begin{longtable}{@{}L{.12\linewidth}L{.79\linewidth}@{}}
\toprule
Identifier & Relationship to the checked source\\\midrule\endhead
PH01--02 & Explicit Clifford matrices, anticommutation, mode self-adjointness, and square identities correspond to the early Clifford modules and Blocks 41 and 44.\\
PH03 & The full spectral-projector formula, rank calculation, and closed trigonometric propagator are ordinary deductions in the manuscript. Individual projector ingredients do not alone certify this full combined statement.\\
PH04--06 & Modewise unitarity, finite Parseval, and packet norm preservation have checked counterparts. Additional finite block deductions follow in the prose.\\
PH07--08 & Positive-mesh quadratic order has checked counterparts. The explicit all-real constants and the displayed elementary proofs are ordinary deductions.\\
PH09 & The full coordinate-exponential conjugacy is proved in the text. The source contains the coordinate ingredients; a matching formal declaration is not asserted here.\\
PH10 & The exact three-mode example is a worked mathematical consequence, accompanied by bounded numerical illustrations.\\
PH11 & The displayed sharp Duhamel coefficient is an ordinary deduction. Block 45 checks compact-time convergence through a different route.\\
PH12 & The expectation calculations and their explicit comparison bound are ordinary finite-dimensional deductions.\\
Later results & Full coefficient/physical evolution, cutoff removal, maximal-domain self-adjointness, graph approximation, and the weighted physical domain correspond to Blocks 46--49.\\
\bottomrule
\end{longtable}

The graph-norm weight comparison and axis-supported domain example are additional explanatory deductions. Their proofs are elementary and displayed in the text. They are not added to the checkpoint's theorem count.

\section{The remaining evolution identification}
The immediate next formal target is to identify the physical unitary family with the evolution generated by $H_m^{\mathrm{phys}}$. In customary notation, the intended conclusion is
\[
\mathcal U(t)=e^{-itH_m^{\mathrm{phys}}},
\]
where the expression on the right has the genuine unbounded spectral or functional-calculus meaning. This displayed equality is a \emph{target}, not a theorem claimed by the present checkpoint.

On a finite Fourier packet, the two constructions suggest the same mode matrices. That agreement is useful input, but a proof must connect the whole Hilbert-space group with its full generator domain. A direct route would establish strong continuity, differentiate on the specified domain with control of the Fourier tail, and prove the appropriate generator equality and uniqueness statement. A functional-calculus route would establish the corresponding spectral representation and compare its action with the constructed modewise group. Either route must use an actual formal theorem rather than treating the familiar notation as its own justification.

This is the right place to concentrate the next block. The operator, its domain, its self-adjointness, and the candidate evolution are already available. The missing theorem ties these independently constructed parts together.

\section{Operator limits and regularity rates}
The checked propagator approximation is a strong convergence statement for each fixed initial field, uniform on compact time intervals. The next operator-level comparison concerns generators. A precise graph or strong-resolvent theorem needs the approximation operators, domains, ambient spaces, and comparison maps to be specified together.

The current centered family acts on the same infinite Fourier index set as the continuum family. It is not a proof about a changing finite sampled grid. If a finite-resolution lattice model is to be used instead, its embedding into the torus Hilbert space and its compatibility with the earlier refinement constructions are part of the theorem's hypotheses or proof obligations.

Quantitative global rates also require more information about the initial field. PH11 exhibits a finite-packet coefficient involving $C(n)^2$, where $C(n)$ grows cubically in momentum. Extending that particular bound directly to an infinite sum would require the corresponding weighted moment to be finite. We make no claim that the first-order Hamiltonian domain alone supplies every quadratic propagator rate. The required regularity depends on the estimate being proved.

Stronger weighted estimates are therefore a separate continuation target. They should state both the precise weight and the norm in which the error is measured. A rate with an unstated regularity assumption would obscure exactly the distinction the maximal-domain construction has clarified.

\section{The broader physics program}
The free model is an important completed analytic sector within a larger collection. Extending it to background gauge coupling requires an operator that incorporates the checked connection and link estimates while retaining a valid domain and evolution theory. Dynamical gauge fields introduce further degrees of freedom, field energy, a Gauss constraint, and matter backreaction. These objects are not supplied by inserting the word gauge into the free Hamiltonian.

The recovered gravity sources develop constraint algebras, local reduction, and controlled strong-field constructions under their own assumptions. The cosmological sources develop a linked chain of model calculations. Those materials belong in their own monograph papers, with the later covariant backreaction correction integrated where it changes the interpretation of an earlier benchmark. The collection guide provides that reconciliation and preserves the older source record.

This volume neither derives a dynamical spacetime metric nor predicts an observational parameter. Its contribution to that wider program is a rigorous free-field foundation and an explicit boundary for the next operator constructions.

\atlasfigure{12_frontier}{The current endpoint and the next obligations. The diagram is a status map, not an estimate of the percentage of physics completed.}

\section{Using the accompanying release}
The PDFs provide four coordinated entry points: the collection guide, this monograph, the visual atlas, and the Lean companion. The source package contains their editable LaTeX, the shared monograph style and fonts, the Python figure generator, numerical figure evidence, the source concordance, and the recovered archives.

The frozen Lean archive is preserved without changes. Some inherited generic wrappers point to older checkpoints, and its aggregate root does not import the Block 48 and Block 49 modules. The version-specific Block 49 handoff and its named module are therefore the authority for replay. The Lean companion gives the exact recorded entrypoint. The collection is not advertised as a newly repaired aggregate Lean build.

Every new figure has a reproducible generator. The numerical checks are scoped to the examples used by those figures. The editorial verification record distinguishes source inspection, document compilation, figure checks, and inherited formal receipts so that the reader can tell which operation produced each piece of evidence.
